Boltzmann Generators, introduced by Noé et al.~\cite{Noe2019Boltzmann}, address the challenges of high-dimensional sampling by combining deep learning techniques with importance sampling. Building upon the Normalizing Flows introduced in \cref{ch:NF}, Boltzmann Generators use a bijective transformation $f_\ta: \eR^n \ra \eR^n$ to map a simple prior distribution $p_Z(\bz)$ (e.g., a multivariate Gaussian distribution) to a distribution $q_\ta(\bx)$ that approximates the target Boltzmann distribution $\rhoeq(\bx) = \frac{1}{Z} e^{-\bt \cU(\bx)}$, where $\cU(\bx)$ is the potential energy function, $Z$ is the partition function, and $\bt$ is the inverse temperature.

The transformed distribution is given by:
\begin{align}
    q_\ta(\bx) = p_Z(f_\ta^{-1}(\bx)) \left|\det \frac{\der f_\ta^{-1}(\bx)}{\der \bx}\right|. \label{eq:bg_transformed_dist}
\end{align}

To train the Boltzmann Generator, we minimize the KL divergence between $q_\ta(\bx)$ and $\rhoeq(\bx)$:
\begin{align}
    \KL[q_\ta||\rhoeq] &= \eE_{q_\ta}[\log q_\ta(\bx) + \bt \cU(\bx)] + \log Z \label{eq:bg_kl}\\
                        &= -\cS[q_\ta] + \beta \eE_{q_\ta}[\cU(\bx)] + \log Z, \label{eq:bg_kl_entropy}
\end{align}
where $\cS[q_\ta]$ is the entropy of $q_\ta$. This objective can be expressed in terms of the latent variables $\bz$:
\begin{align}
    \KL[q_\ta || \rhoeq] = -\eE_{p_Z(\bz)}[\log |\det J_{f_\ta}(\bz)|] + \beta \eE_{p_Z(\bz)}[\cU(f_\ta(\bz))] + \log Z, \label{eq:bg_kl_latent}
\end{align}
where $J_{f_\ta}(\bz)$ is the Jacobian of $f_\ta$ at $\bz$.

The gradients with respect to $\ta$ can be computed using the change of variables formula and the Jacobian of the normalizing flow:
\begin{align}
    \na_\ta \KL[q_\ta || \rhoeq] = -\eE_{p_Z(\bz)}[\na_\ta \log |\det J_{f_\ta}(\bz)|] + \beta \eE_{p_Z(\bz)}[\na_\ta \cU(f_\ta(\bz))]. \label{eq:bg_kl_gradient}
\end{align}
Note that the $\log Z$ term is constant with respect to $\ta$ and thus disappears in the gradient computation.

To compute these gradients, we can use the following steps:
\begin{align}
    \na_\ta \log |\det J_{f_\ta}(\bz)| &= \tr\left(J_{f_\ta}(\bz)^{-1} \na_\ta J_{f_\ta}(\bz)\right), \label{eq:bg_jacobian_grad}\\
    \na_\ta \cU(f_\ta(\bz)) &= \na_\bx \cU(f_\ta(\bz)) \cdot \na_\ta f_\ta(\bz), \label{eq:bg_potential_grad}
\end{align}
where $\tr(\cdot)$ denotes the trace of a matrix. The term $\na_\ta J_{f_\ta}(\bz)$ can be computed using automatic differentiation techniques.

When samples $\{\bx_i\}_{i=1}^M$ from the target distribution $\rhoeq(\bx)$ are available (e.g., from MCMC, molecular dynamics), we can train the Boltzmann Generator by minimizing the reverse KL divergence:
\begin{align}
    \KL[\rhoeq || q_\ta] &= \eE_{\rhoeq}[\log \rhoeq(\bx) - \log q_\ta(\bx)] \label{eq:bg_reverse_kl}\\
                            &= -\frac{1}{M} \sum_{i=1}^M \log q_\ta(\bx_i) + \text{const}, \label{eq:bg_reverse_kl_samples}
\end{align}
where the constant term includes $\eE_{\rhoeq}[\log \rhoeq(\bx)]$, which is independent of $\ta$. This is equivalent to maximizing the log-likelihood:
\begin{align}
    \cL(\ta) &= \frac{1}{M} \sum_{i=1}^M \log q_\ta(\bx_i) \label{eq:bg_mle}\\
                &= \frac{1}{M} \sum_{i=1}^M \log p_Z(f_\ta^{-1}(\bx_i)) + \log |\det J_{f_\ta^{-1}}(\bx_i)|. \label{eq:bg_mle_expanded}
\end{align}

The gradients for this objective are given by:
\begin{align}
    \na_\ta \cL(\ta) = \frac{1}{M} \sum_{i=1}^M \na_\ta \log p_Z(f_\ta^{-1}(\bx_i)) + \na_\ta \log |\det J_{f_\ta^{-1}}(\bx_i)|. \label{eq:bg_mle_gradient}
\end{align}

Once trained, the Boltzmann Generator can be used for importance sampling. Given samples $\{\bx_i\}_{i=1}^N$ drawn from $q_\ta(\bx)$, expectations under $\rhoeq(\bx)$ are estimated using:
\begin{align}
    \eE_{\rhoeq}[h(\bx)] \approx \frac{\sum_{i=1}^N w(\bx_i) h(\bx_i)}{\sum_{i=1}^N w(\bx_i)}, \label{eq:bg_importance_sampling}
\end{align}
where $w(\bx_i) = \frac{e^{-\bt \cU(\bx_i)}}{q_\ta(\bx_i)}$ are the importance weights.

The partition function $Z$ can be estimated using importance sampling:
\begin{align}
    Z \approx \frac{1}{N} \sum_{i=1}^N \frac{e^{-\bt \cU(\bx_i)}}{q_\ta(\bx_i)}, \quad \bx_i \sim q_\ta(\bx). \label{eq:bg_partition_function}
\end{align}
This estimation of $Z$ is performed after the Boltzmann Generator has been trained and is separate from the training process.